\documentclass[a4paper,12pt,titlepage]{report}

% ============================================================
% Компилировать XeLaTeX или LuaLaTeX.
% ============================================================
\usepackage{iftex}
\ifPDFTeX
  \PackageError{main}{Use XeLaTeX or LuaLaTeX instead of pdfLaTeX}
    {Compile the document with XeLaTeX or LuaLaTeX.}
\fi

% ============================================================
% Шрифт и языки
% ============================================================
\usepackage{fontspec}
\defaultfontfeatures{Ligatures=TeX}
\setmainfont{Times New Roman}

\usepackage[main=russian,english]{babel}

% ============================================================
% Основные пакеты
% ============================================================
\usepackage{amsmath,amssymb,amsfonts,amsthm}
\usepackage{mathtext}
\usepackage{enumerate}
\usepackage{enumitem}
\usepackage{float}
\usepackage{graphicx}
\usepackage{placeins}
\usepackage{geometry}
\usepackage[onehalfspacing]{setspace}
\usepackage{indentfirst}
\usepackage{titlesec}
\usepackage[font=small,tableposition=top]{caption}
\usepackage{subcaption}
\usepackage[title,titletoc]{appendix}
\usepackage[square,numbers,sort&compress]{natbib}
\usepackage{listings}
\usepackage{xcolor}

% Путь к рисункам.
\graphicspath{{figure/}}

% ============================================================
% Поля страницы
% ============================================================
\geometry{
  left=2.5cm,
  right=1.5cm,
  top=2cm,
  bottom=2cm
}

% ============================================================
% Форматирование заголовков
% ============================================================
\titleformat{\chapter}[display]
  {\filcenter\bfseries}
  {\MakeUppercase{\chaptertitlename} \thechapter}
  {8pt}
  {\MakeUppercase}

\titleformat{\section}
  {\normalsize\bfseries}
  {\thesection.}
  {1em}
  {}

\titleformat{\subsection}
  {\normalsize\bfseries}
  {\thesubsection.}
  {1em}
  {}

\titlespacing*{\chapter}{0pt}{-30pt}{8pt}
\titlespacing*{\section}{\parindent}{4\baselineskip}{\baselineskip}
\titlespacing*{\subsection}{\parindent}{4\baselineskip}{\baselineskip}

% ============================================================
% Оглавление
% ============================================================
\renewcommand{\contentsname}{Содержание}

% В оглавлении выводить главы и разделы.
\setcounter{tocdepth}{1}

% ============================================================
% Подписи рисунков и таблиц
% ============================================================
\DeclareCaptionLabelFormat{gostfigure}{Рисунок #2}
\DeclareCaptionLabelFormat{gosttable}{Таблица #2}
\DeclareCaptionLabelSeparator{gost}{. }

\captionsetup{labelsep=gost}
\captionsetup[figure]{labelformat=gostfigure}
\captionsetup[table]{labelformat=gosttable}

\renewcommand{\thesubfigure}{\asbuk{subfigure}}

% ============================================================
% Списки
% ============================================================
\setlist{nolistsep}
\renewcommand{\labelitemi}{-}
\renewcommand{\labelenumi}{\asbuk{enumi})}
\renewcommand{\labelenumii}{\arabic{enumii})}

% ============================================================
% Приложения
% ============================================================
\newcommand{\append}[1]{%
  \clearpage
  \chapter*{\MakeUppercase{#1}}%
  \addcontentsline{toc}{chapter}{\MakeUppercase{#1}}%
}

% ============================================================
% Библиография
% ============================================================
\bibliographystyle{utf8gost705u}
\renewcommand{\bibnumfmt}[1]{#1.\hfill}
\setlength{\bibsep}{0pt}

\renewcommand{\bibsection}{%
  \chapter*{Список литературы}%
  \addcontentsline{toc}{chapter}{Список литературы}%
}

% ============================================================
% Листинги кода
% ============================================================
\lstdefinestyle{customc}{
  belowcaptionskip=1\baselineskip,
  breaklines=true,
  frame=L,
  xleftmargin=\parindent,
  showstringspaces=false,
  basicstyle=\footnotesize\ttfamily,
  keywordstyle=\bfseries\color{green!40!black},
  commentstyle=\itshape\color{purple!40!black},
  identifierstyle=\color{blue},
  stringstyle=\color{orange},
  numbers=left,
  numberstyle=\tiny
}

\lstdefinestyle{customasm}{
  belowcaptionskip=1\baselineskip,
  breaklines=true,
  frame=L,
  xleftmargin=\parindent,
  language=[x86masm]Assembler,
  basicstyle=\footnotesize\ttfamily,
  commentstyle=\itshape\color{purple!40!black}
}

% ============================================================
% Теоремы и формулы
% ============================================================
\theoremstyle{plain}

\newtheorem{thm}{Theorem}[section]
\newtheorem{St}{Утверждение}
\newtheorem{Th}{Теорема}
\newtheorem{Cor}{Следствие}
\newtheorem{Lm}{Лемма}

% Сквозная нумерация формул.
\counterwithout{equation}{section}
\counterwithout{equation}{chapter}

% ============================================================
% Начало документа
% ============================================================
\begin{document}

\singlespacing

% Если файл называется title.tex:
\begin{titlepage}

\begin{center}

% Логотип временно выключен, чтобы сборка не зависела
% от формата изображения. Для XeLaTeX используйте PDF или PNG:
%
\includegraphics[width=0.8cm]{logo_unn}\par
\vspace{0.5em}

МИНИСТЕРСТВО НАУКИ И ВЫСШЕГО ОБРАЗОВАНИЯ\\
РОССИЙСКОЙ ФЕДЕРАЦИИ

\vspace{1em}

Федеральное государственное автономное образовательное учреждение\\
высшего образования

\vspace{0.5em}

\textbf{<<Национальный исследовательский\\
Нижегородский государственный университет им. Н.\,И.~Лобачевского>>}

\vspace{2em}

\textbf{Институт информационных технологий, математики и механики}

\vspace{2em}

Направление подготовки: <<Прикладная математика\\
и информатика>>

\end{center}

\vspace{1em}

\begin{center}
{\Large\bfseries
Отчёт по лабораторной работе №\,4\\
по уравнениям математической физики
}
\end{center}

\vspace{2.5em}

\begin{center}
{\large\bfseries Тема:}

\vspace{0.5em}

{\large\bfseries
<<Распространение тепла в конечном стержне>>
}
\end{center}

\vspace{4em}

\hfill
\begin{minipage}[t]{0.42\linewidth}
\raggedright

Выполнили:\\
студенты группы 3823Б1ПМоп2\\
Золкин Иван Александрович\\
Носов Владимир Михайлович

\vspace{1.5em}

Проверил:\\
к.\,ф.-м.\,н., доцент\\
Муняев Вячеслав Олегович

\end{minipage}

\vfill

\begin{center}
Нижний Новгород\\
2026
\end{center}

\end{titlepage}

% Если у вас файл называется title-2.tex, используйте вместо
% предыдущей строки следующую:
% \input{title-2}

% Титульный лист — первая страница, номер не показывается.
\setcounter{page}{2}

\onehalfspacing
\setlength{\parindent}{1.25cm}

\tableofcontents

\chapter{Уравнение теплопроводности и индивидуальные границы}
\section{Тепловой баланс и закон Ньютона}
Пусть $S$ --- площадь поперечного сечения, $P$ --- его периметр,
$k>0$ --- теплопроводность, $c>0$ --- объёмная теплоёмкость,
$H\ge0$ --- коэффициент теплоотдачи, $F(x,t)$ --- объёмная мощность
источников, $T_e(x,t)$ --- температура среды. Для тонкого однородного
стержня температура $u$ приближённо постоянна по сечению.
Закон Фурье даёт поток вдоль оси $Q_x=-kS u_x$.
Приток через торцы малого элемента равен $kS u_{xx}\,dx$;
боковая потеря по закону Ньютона равна $HP(u-T_e)\,dx$.
Накопление тепла равно $cS u_t\,dx$. Поэтому
\[
 cS u_t=kS u_{xx}-HP(u-T_e)+FS,
\qquad u_t=a^2u_{xx}-bu+f,
\]
где
\[
 a^2=\frac{k}{c},\quad b=\frac{HP}{cS},\quad f=\frac Fc+bT_e.
\]
Для круглого сечения $S=\pi R^2$, $P=2\pi R$,
$b=2H/(cR)$. Это соответствует лекции 22, формуле (18),
с заменой лекционного $c_{\rm уд}\rho$ на $c$.
При теплоизоляции боковой поверхности $H=b=0$;
при отсутствии источников и бокового теплообмена также $f=0$.
\section{Условия варианта 11}
В столбцах 1--2 строки 11 указаны первый и второй род соответственно.
Поэтому произвольные индивидуальные данные записываем как
\[
 \boxed{u(0,t)=\mu(t),\qquad -u_x(l,t)=\nu(t),\qquad u(x,0)=\phi(x).}
\]
Знак справа сохраняет нормировку задания
$\gamma u(l,t)-\delta u_x(l,t)=\nu(t)$ при $\gamma=0$, $\delta=1$.
Это не означает, что $\nu$ уже является физическим потоком:
наружная плотность потока на правом торце равна $-ku_x(l,t)=k\nu(t)$.
При ином соглашении $u_x(l,t)=q(t)$ надо положить $\nu=-q$.
Для практического пункта 3 строка варианта задаёт
\[
 \mu(t)=A\cos(\omega_1t)+B\sin(\omega_2t),\qquad \nu(t)=0.
\]
В пунктах 2 и 4 по отдельному требованию задания вместо этих функций
берутся постоянные неоднородности $\mu=A$, $\nu=B$.
\section{Замечание о третьем роде}
Физический закон Ньютона на границе имеет вид
$k\partial_nu+H(u-T_e)=0$, где $n$ --- внешняя нормаль.
Следовательно, слева это $-ku_x+H(u-T_e)=0$,
справа --- $ku_x+H(u-T_e)=0$.
Запись с минусом перед $u_x$ на обоих концах и положительными
коэффициентами не является общей физической записью охлаждения обоих
торцов. В варианте 11 третьего рода нет, поэтому эту неоднозначность
не переносим в индивидуальное решение и не предполагаем положительность
спектра произвольной задачи с такими знаками.
\section{Согласование данных}
Для классического решения вплоть до $t=0$ необходимы
$\phi(0)=\mu(0)$ и $-\phi'(l)=\nu(0)$.
Без согласования рассматриваем решение при $t>0$, с начальным условием
в $L^2(0,l)$; непрерывность всех производных в угловых точках не заявляется.

\chapter{Метод Фурье и представление через функцию Грина}
\section{Устранение неоднородности границ}
Для варианта 11 решаем
\[
 u_t=a^2u_{xx}-bu+f,\quad u(0,t)=\mu(t),\quad -u_x(l,t)=\nu(t),
 \quad u(x,0)=\phi(x),\qquad a>0,\ b\ge0.
\]
Пусть $\mu,\nu$ достаточно регулярны (например, $C^1$ по времени),
а $\phi\in L^2(0,l)$ и источник допускает интегральное представление ниже.
Линейная функция
\[
 \boxed{w(x,t)=\mu(t)-x\nu(t)}
\]
удовлетворяет обеим неоднородным границам. При $u=w+v$ получаем
\[
 v_t=a^2v_{xx}-bv+\widetilde f,\quad v(0,t)=0,\quad v_x(l,t)=0,
 \quad v(x,0)=\widetilde\phi(x),
\]
\[
 \widetilde f=f-\mu'-b\mu+x(\nu'+b\nu),
 \qquad \widetilde\phi=\phi-\mu(0)+x\nu(0).
\]
Такой переход является индивидуальным аналогом замены из лекции 24.
\section{Спектральная задача}
Ищем $-X''=\lambda X$, $X(0)=0$, $X'(l)=0$.
Интегрирование по частям даёт
$\lambda\int_0^lX^2dx=\int_0^l(X')^2dx>0$
для ненулевой собственной функции. При $\lambda=k^2$
имеем $X=C\sin(kx)$ и $\cos(kl)=0$. Поэтому
\[
 k_n=\frac{(n+\tfrac12)\pi}{l},\quad \lambda_n=k_n^2,
 \quad X_n(x)=\sin(k_nx),\quad N_n=\frac l2,\quad n=0,1,\ldots.
\]
Здесь нет ни нулевой, ни отрицательной моды.
\section{Модальные уравнения}
Обозначим
\[
 \sigma_n=a^2k_n^2+b,\qquad
 c_n=\frac2l\int_0^l\widetilde\phi(\xi)\sin(k_n\xi)\,d\xi,
\quad F_n(t)=\frac2l\int_0^l\widetilde f(\xi,t)\sin(k_n\xi)\,d\xi.
\]
Разложение $v=\sum_nT_nX_n$ даёт
\[
 T_n'+\sigma_nT_n=F_n,\quad T_n(0)=c_n,
\qquad T_n(t)=c_ne^{-\sigma_nt}+\int_0^te^{-\sigma_n(t-\tau)}F_n(\tau)\,d\tau.
\]
Следовательно,
\[
 \boxed{u(x,t)=\mu(t)-x\nu(t)+\sum_{n=0}^{\infty}\sin(k_nx)
 \left[c_ne^{-\sigma_nt}+\int_0^te^{-\sigma_n(t-\tau)}F_n(\tau)\,d\tau\right].}
\]
Если $f_n=(2/l)\int_0^lf(\xi,t)\sin(k_n\xi)d\xi$, то
\[
 F_n=f_n-\frac{2(\mu'+b\mu)}{lk_n}
       +\frac{2(-1)^n(\nu'+b\nu)}{lk_n^2}.
\]
Использованы интегралы
$\int_0^l\sin(k_n\xi)d\xi=1/k_n$,
$\int_0^l\xi\sin(k_n\xi)d\xi=(-1)^n/k_n^2$.
\section{Функция Грина}
Введём тепловое ядро
\[
 \boxed{G(x,\xi,t)=\frac2l\sum_{n=0}^{\infty}
 \sin(k_nx)\sin(k_n\xi)e^{-(a^2k_n^2+b)t},\quad t>0.}
\]
Тогда окончательная формула имеет вид
\[
\boxed{\begin{aligned}
 u(x,t)={}&\mu(t)-x\nu(t)
 +\int_0^lG(x,\xi,t)[\phi(\xi)-\mu(0)+\xi\nu(0)]\,d\xi\\
 &+\int_0^t\int_0^lG(x,\xi,t-\tau)
 [f(\xi,\tau)-\mu'(\tau)-b\mu(\tau)
       +\xi(\nu'(\tau)+b\nu(\tau))]\,d\xi\,d\tau.
\end{aligned}}
\]
Начальные данные восстанавливаются в $L^2$, а при достаточной гладкости
и согласовании --- в классическом смысле. Формула учитывает произвольные
индивидуальные неоднородности; подставлять гармоники нужно лишь при
решении соответствующего практического пункта.

\chapter{Физический смысл и график функции Грина}
\section{Температурный импульс}
Функция $G(x,\xi,t)$ --- отклик на начальное распределение
$\delta(x-\xi)$ при однородных границах. Она удовлетворяет
$G_t=a^2G_{xx}-bG$, $G(0,\xi,t)=0$, $G_x(l,\xi,t)=0$ и
\[
 \lim_{t\downarrow0}\int_0^lG(x,\xi,t)\psi(x)dx=\psi(\xi),\qquad 0<\xi<l.
\]
Единичный температурный импульс не равен одному джоулю:
внесение теплоты $Q$ при объёмной теплоёмкости $c$ и сечении $S$
даёт $u=QG/(cS)$. Источник в момент $\tau$ порождает ядро
$G(x,\xi,t-\tau)$. Линейность позволяет суммировать такие отклики.
\section{Приближённое построение}
Используем первые $N=100$ мод при $l=a^2=1$, $b=0.5$, $\xi=0.4$,
$t_1=0.01$, $t_2=0.06$. Это безразмерные параметры численного примера,
а не заданные физические константы варианта.
График вычислен по спектральной формуле; внешняя картинка не требуется.
Для графика взята 81 точка $x$, а численные характеристики ниже проверены
на более густой сетке из 2001 точки.
\begin{figure}[htbp]
\centering
\begin{tikzpicture}
\begin{axis}[width=.9\textwidth,height=7cm,xlabel={$x$},ylabel={$G_N$},grid=major,legend pos=north east]
\addplot[blue,thick] coordinates {
(0.00000,0) (0.01250,0.0258722216228) (0.02500,0.0527481184805) (0.03750,0.0816294039483) (0.05000,0.113512707456) (0.06250,0.149384200598) (0.07500,0.190211016216) (0.08750,0.236928707166) (0.10000,0.290424254848) (0.11250,0.351514455725) (0.12500,0.420919876742) (0.13750,0.499234963358) (0.15000,0.586895288126) (0.16250,0.684143320111) (0.17500,0.790994449368) (0.18750,0.907205287304) (0.20000,1.03224645404) (0.21250,1.16528213124) (0.22500,1.3051585819) (0.23750,1.45040360341) (0.25000,1.59923848416) (0.26250,1.74960348573) (0.27500,1.8991971945) (0.28750,2.04552931422) (0.30000,2.1859856501) (0.31250,2.31790322073) (0.32500,2.4386526866) (0.33750,2.54572466095) (0.35000,2.63681602453) (0.36250,2.70991214396) (0.37500,2.76336092018) (0.38750,2.7959348808) (0.40000,2.80687806543) (0.41250,2.79593520871) (0.42500,2.7633616511) (0.43750,2.70991344265) (0.45000,2.63681817697) (0.46250,2.54572812753) (0.47500,2.43865818511) (0.48750,2.3179118509) (0.50000,2.18599907684) (0.51250,2.04555003328) (0.52500,1.89922891341) (0.53750,1.74965166412) (0.55000,1.59931109233) (0.56250,1.45051217666) (0.57500,1.30531967102) (0.58750,1.1655192775) (0.60000,1.03259285035) (0.61250,0.907707326091) (0.62500,0.791716401399) (0.63750,0.685173436949) (0.65000,0.588353671512) (0.66250,0.501283595639) (0.67500,0.423775253368) (0.68750,0.355463298809) (0.70000,0.295842809334) (0.71250,0.244306120395) (0.72500,0.200177268554) (0.73750,0.162742979426) (0.75000,0.131279489008) (0.76250,0.105074817716) (0.77500,0.0834464091688) (0.78750,0.0657542889959) (0.80000,0.0514100867091) (0.81250,0.0398823952891) (0.82500,0.0306990216056) (0.83750,0.0234467124599) (0.85000,0.0177689339833) (0.86250,0.013362245401) (0.87500,0.00997175084577) (0.88750,0.00738604340571) (0.90000,0.00543198122363) (0.91250,0.00396956214087) (0.92500,0.00288709554009) (0.93750,0.0020968106724) (0.95000,0.00153099155409) (0.96250,0.00113869008945) (0.97500,0.000883041153044) (0.98750,0.000739185041154) (1.00000,0.000692792622264)
};
\addlegendentry{$t=0.01$}
\addplot[red,thick] coordinates {
(0.00000,0) (0.01250,0.0478081373586) (0.02500,0.0955127838547) (0.03750,0.14301073166) (0.05000,0.190199339055) (0.06250,0.236976813562) (0.07500,0.283242495119) (0.08750,0.328897139215) (0.10000,0.37384319985) (0.11250,0.417985112051) (0.12500,0.461229573639) (0.13750,0.503485825763) (0.15000,0.544665931633) (0.16250,0.584685052749) (0.17500,0.623461721794) (0.18750,0.660918111224) (0.20000,0.696980296466) (0.21250,0.731578512529) (0.22500,0.764647402712) (0.23750,0.796126257997) (0.25000,0.825959245651) (0.26250,0.854095625493) (0.27500,0.880489952226) (0.28750,0.90510226225) (0.30000,0.927898243338) (0.31250,0.948849385613) (0.32500,0.967933112302) (0.33750,0.985132888836) (0.35000,1.00043830894) (0.36250,1.01384515651) (0.37500,1.0253554422) (0.38750,1.0349774138) (0.40000,1.04272553978) (0.41250,1.04862046528) (0.42500,1.05268894046) (0.43750,1.05496372103) (0.45000,1.05548344101) (0.46250,1.05429245827) (0.47500,1.05144067334) (0.48750,1.0469833223) (0.50000,1.04098074499) (0.51250,1.03349812945) (0.52500,1.02460523447) (0.53750,1.01437609155) (0.55000,1.00288868826) (0.56250,0.990224634868) (0.57500,0.976468816335) (0.58750,0.961709031782) (0.60000,0.946035623719) (0.61250,0.929541099254) (0.62500,0.912319745579) (0.63750,0.894467242022) (0.65000,0.876080270877) (0.66250,0.857256129226) (0.67500,0.838092343816) (0.68750,0.818686291025) (0.70000,0.799134823752) (0.71250,0.779533907006) (0.72500,0.759978263743) (0.73750,0.740561032403) (0.75000,0.721373437367) (0.76250,0.702504473438) (0.77500,0.684040605214) (0.78750,0.666065482087) (0.80000,0.648659669409) (0.81250,0.631900396202) (0.82500,0.615861319628) (0.83750,0.600612306302) (0.85000,0.586219230384) (0.86250,0.572743788278) (0.87500,0.560243329659) (0.88750,0.548770704483) (0.90000,0.538374125552) (0.91250,0.529097046172) (0.92500,0.520978052418) (0.93750,0.514050769516) (0.95000,0.50834378185) (0.96250,0.503880566144) (0.97500,0.500679437408) (0.98750,0.498753507282) (1.00000,0.498110654507)
};
\addlegendentry{$t=0.06$}
\end{axis}\end{tikzpicture}
\caption{Отклик на температурный импульс при двух моментах времени.}
\end{figure}
\section{Интерпретация и погрешность}
Максимумы равны приблизительно $2.806878$ и $1.055518$;
их координаты на сетке --- $0.4000$ и $0.4475$.
Диффузия расширяет возмущение, а охлаждение и потери на левом торце
уменьшают его амплитуду. Правый торец теплоизолирован.
Смещение пика обусловлено несимметричными граничными условиями.
Интегральная температура $M=\int_0^lGdx$ подчиняется балансу
\[
 M'=-a^2G_x(0,\xi,t)-bM.
\]
Метод трапеций на 2001 точке даёт $M(t_1)=0.990358$,
$M(t_2)=0.729564$. Проверка усечения отделена от ошибки сетки:
\[
 \sup_x|G-G_N|\le\frac2l e^{-bt}\sum_{n=N}^{\infty}e^{-a^2k_n^2t}.
\]
Сравнение 100 и 200 мод при этих временах не выявляет различий
в машинной точности. При $t\downarrow0$ число мод надо увеличивать.

\chapter{Стационарный режим теплопроводности}

\section{Определение стационарного режима}

Стационарным режимом называется распределение температуры, не зависящее
от времени:
\[
 u(x,t)=U(x),\qquad u_t(x,t)=0.
\]
Температура в каждой точке стержня остаётся постоянной, однако это
не означает отсутствия тепловых потоков. Потоки, теплоотдача и источники
тепла взаимно компенсируются, так что накопления тепла не происходит.

Для задачи
\[
\begin{cases}
 u_t=a^2u_{xx}-bu+f(x,t),&0<x<l,\\
 u(x,0)=\phi(x),\\
 \alpha u(0,t)-\beta u_x(0,t)=\mu(t),\\
 \gamma u(l,t)-\delta u_x(l,t)=\nu(t)
\end{cases}
\]
точный стационарный режим возможен лишь при независимых от времени
правых частях:
\[
 f(x,t)=f_s(x),\qquad\mu(t)=\mu_s,\qquad\nu(t)=\nu_s.
\]
Действительно, подстановка $u=U(x)$ делает левые части уравнения
и граничных условий независимыми от времени.
Профиль $U$ определяется краевой задачей
\[
\boxed{\begin{cases}
 a^2U''(x)-bU(x)+f_s(x)=0,&0<x<l,\\
 \alpha U(0)-\beta U'(0)=\mu_s,\\
 \gamma U(l)-\delta U'(l)=\nu_s.
\end{cases}}
\]
Начальное условие в эту краевую задачу не входит.
Если $\phi=U$, то решение стационарно с начального момента;
если $\phi\ne U$, возникает переходный процесс.

Следует различать существование стационарного решения и установление
стационарного режима:
\[
 \lim_{t\to\infty}u(\cdot,t)=U.
\]
Для второго утверждения необходима устойчивость стационарного решения.
Если данные лишь стремятся к $f_s,\mu_s,\nu_s$, то предельный профиль
определяется той же краевой задачей; сходимость требует дополнительных
условий устойчивости и регулярности данных.

\section{Влияние неоднородности уравнения}

Функция $f_s(x)$ описывает распределённое тепловое воздействие:
внутренние источники и, при принятом ранее переобозначении,
вклад температуры окружающей среды. Она непосредственно определяет
форму стационарного профиля:
\[
 U''(x)=\frac{b}{a^2}U(x)-\frac{f_s(x)}{a^2}.
\]
При $b=0$ имеем
\[
 U''(x)=-\frac{f_s(x)}{a^2}.
\]
Поэтому при $f_s=0$ профиль линейный, а при постоянном ненулевом
источнике --- квадратичный. При $b>0$ даже при $f_s=0$ профиль
в общем случае является комбинацией гиперболических функций.

Для условий первого рода при $x=0$ и второго рода при $x=l$ положим
\[
 U(0)=T_0,\qquad U'(l)=q_0.
\]
Здесь $q_0$ --- заданная производная температуры, а не непосредственно
тепловой поток. Поток в положительном направлении оси равен
$-kS q_0$, где $k$ --- теплопроводность и $S$ --- площадь сечения.
При $b=0$ интегрирование даёт
\[
\boxed{U(x)=T_0+q_0x+
 \frac{1}{a^2}\int_0^l\min(x,\xi)f_s(\xi)\,d\xi.}
\]
В частности, если $f_s(x)=F_0$, то
\[
 U(x)=T_0+q_0x+\frac{F_0}{a^2}
 \left(lx-\frac{x^2}{2}\right).
\]
Таким образом, изменение источника изменяет стационарный профиль
даже при неизменных данных на границах.

\section{Влияние начальных условий}

Пусть стационарный профиль $U$ существует. Выполним замену
\[
 u(x,t)=U(x)+v(x,t).
\]
При постоянных во времени внешних данных функция $v$ удовлетворяет
\[
 v_t=a^2v_{xx}-bv,\qquad v(x,0)=\phi(x)-U(x)
\]
с однородными граничными условиями.
Пусть $X_n$ --- собственные функции задачи
\[
 -X_n''=\lambda_nX_n,\qquad
 \alpha X_n(0)-\beta X_n'(0)=0,\qquad
 \gamma X_n(l)-\delta X_n'(l)=0.
\]
Тогда
\[
 v(x,t)=\sum_n c_n e^{-(a^2\lambda_n+b)t}X_n(x),
 \qquad
 c_n=\frac{\int_0^l(\phi-U)X_n\,dx}{\int_0^lX_n^2\,dx}.
\]
Если для всех мод выполняется
\[
 a^2\lambda_n+b>0,
\]
переходная часть затухает. В этом случае стационарный профиль единственен
и не зависит от начальных условий; начальная температура определяет
амплитуды переходных мод и время достижения заданной близости к равновесию.

Для рассматриваемых условий первого--второго рода
\[
 \lambda_n=\left(\frac{(n+\tfrac12)\pi}{l}\right)^2,
 \qquad n=0,1,2,\ldots.
\]
При $b\ge0$ все моды затухают. В норме $L^2(0,l)$
\[
 \|u(\cdot,t)-U\|_{L^2}
 \le e^{-\sigma_0t}\|\phi-U\|_{L^2},\qquad
 \sigma_0=b+\frac{a^2\pi^2}{4l^2}>0.
\]

Утверждение о независимости предельного режима от начальных данных
не является универсальным. Если оба конца теплоизолированы,
$b=0$ и $f_s=0$, то средняя температура сохраняется и
\[
 u(x,t)\longrightarrow\frac1l\int_0^l\phi(\xi)\,d\xi.
\]
Здесь существует семейство постоянных стационарных решений, а начальные
условия выбирают одно из них. Если имеется незатухающая или растущая мода,
общая сходимость к единственному стационарному режиму отсутствует.
В частности, для записи с минусом перед производной на обоих концах
не следует предполагать положительность спектра без проверки.

\section{Влияние граничных неоднородностей}

Постоянные значения $\mu_s$ и $\nu_s$ непосредственно входят в
стационарную краевую задачу, поэтому определяют уровень температуры
и потоки на концах стержня. Даже если $f_s=0$, неоднородные граничные
условия могут поддерживать ненулевой стационарный профиль.

При выбранных коэффициентах
\[
 \alpha>0,\quad\beta=0,\qquad\gamma=0,\quad\delta>0
\]
получаем
\[
 T_0=\frac{\mu_s}{\alpha},\qquad
 q_0=-\frac{\nu_s}{\delta}.
\]
При $f_s=0$ и $b=0$
\[
\boxed{U(x)=\frac{\mu_s}{\alpha}-\frac{\nu_s}{\delta}x.}
\]
При $f_s=0$ и $b>0$, обозначая $\kappa=\sqrt{b}/a$, имеем
\[
\boxed{U(x)=T_0\cosh(\kappa x)+
 \frac{q_0/\kappa-T_0\sinh(\kappa l)}{\cosh(\kappa l)}
 \sinh(\kappa x).}
\]
Эта формула показывает совместное влияние граничных данных
и теплоотдачи через боковую поверхность.

При переменных во времени $\mu(t)$ и $\nu(t)$ точный стационарный
профиль невозможен. Если они стремятся к постоянным значениям
и система устойчива, возможно установление предельного стационарного
режима. Периодическое граничное воздействие, напротив, может приводить
к установившемуся периодическому режиму, который не является стационарным.

\section{Баланс и существование режима}

Интегрирование стационарного уравнения даёт необходимый баланс:
\[
 a^2\bigl(U'(l)-U'(0)\bigr)-b\int_0^lU(x)\,dx+
 \int_0^lf_s(x)\,dx=0.
\]
Это равенство выражает отсутствие накопления тепла:
поступление через источники и границы компенсируется потерями.
Оно не заменяет полного решения краевой задачи.

При теплоизоляции обоих концов и $b=0$ необходимо
\[
 \int_0^lf_s(x)\,dx=0.
\]
Если это условие нарушено, средняя температура изменяется со временем,
и стационарного режима нет. При выполнении условия профиль определяется
с точностью до аддитивной константы, выбираемой начальной средней
температурой.

Для условий первого рода слева и второго рода справа при $b\ge0$
стационарная задача имеет единственное решение для достаточно регулярных
$f_s$ и любых постоянных граничных данных. Этот профиль определяется
источниками, граничными условиями и коэффициентами модели;
начальные данные влияют только на переходный процесс.

\chapter{Ненулевой стационарный режим изолированного стержня}

\section{Постановка задачи и выбор границ}

Пусть стержень занимает отрезок $[0,l]$, где $l>0$, и имеет
теплоизолированную боковую поверхность. Внешние и внутренние источники
тепла отсутствуют. Поэтому в ранее полученном уравнении необходимо
положить $b=0$ и $f=0$:
\[
 u_t=a^2u_{xx},\qquad 0<x<l,\quad t>0,\qquad a>0.
\]
Начальная температура задана условием
\[
 u(x,0)=x.
\]

В первоначальном варианте задания были указаны условия первого рода
при $x=0$ и второго рода при $x=l$. При их однородности получаем
\[
 u(0,t)=0,\qquad u_x(l,t)=0.
\]
Такие условия не дают ненулевого стационарного режима: из
$U''=0$ следует $U=Ax+B$, а условия $U(0)=0$ и $U'(l)=0$
дают $A=B=0$.

Следовательно, выполнить новый пункт с сохранением прежних типов
границ невозможно. Поскольку в этом пункте требуется подобрать
однородные условия, выберем теплоизоляцию обоих торцов:
\[
\boxed{u_x(0,t)=0,\qquad u_x(l,t)=0.}
\]
По закону Фурье нулевая производная температуры означает отсутствие
теплового потока через торец. В обозначениях общей задачи с изображения
этот выбор соответствует
\[
 \alpha=\gamma=0,\qquad\beta=\delta=1,
 \qquad\mu=\nu=0.
\]
Получаем задачу
\[
\boxed{\begin{cases}
 u_t=a^2u_{xx},&0<x<l,\ t>0,\\
 u_x(0,t)=u_x(l,t)=0,&t>0,\\
 u(x,0)=x,&0\le x\le l.
\end{cases}}
\]

\section{Стационарные решения и сохранение тепла}

Стационарное решение не зависит от времени: $u(x,t)=U(x)$.
Для него
\[
 U''=0,\qquad U'(0)=U'(l)=0.
\]
Поэтому
\[
\boxed{U(x)=C,\qquad C\in\mathbb R.}
\]
Таким образом, стационарная краевая задача имеет семейство решений,
включающее ненулевые постоянные. Значение $C$ выбирается начальными данными.

Для произвольной начальной температуры $\phi(x)$ интегрируем уравнение
по длине стержня:
\[
 \frac{d}{dt}\int_0^lu(x,t)\,dx
 =a^2\bigl(u_x(l,t)-u_x(0,t)\bigr)=0.
\]
Следовательно,
\[
 \int_0^lu(x,t)\,dx=\int_0^l\phi(x)\,dx.
\]
При постоянной объёмной теплоёмкости $c$ и площади сечения $S$
величина $cS\int_0^lu\,dx$ выражает тепловую энергию относительно
выбранного температурного отсчёта. В полностью изолированном стержне
она сохраняется.

Если $u(\cdot,t)$ стремится к постоянной $C$, сохранение средней
температуры даёт
\[
 Cl=\int_0^l\phi(x)\,dx,
 \qquad
\boxed{C=\overline\phi=\frac1l\int_0^l\phi(x)\,dx.}
\]
Для $\phi(x)=x$ получаем
\[
\boxed{C=\frac1l\int_0^lx\,dx=\frac l2\ne0.}
\]

\section{Решение методом Фурье}

Рассмотрим задачу на собственные значения
\[
 -X''=\lambda X,\qquad X'(0)=X'(l)=0.
\]
Она имеет нулевую моду
\[
 \lambda_0=0,\qquad X_0(x)=1
\]
и положительные моды
\[
 \lambda_n=\left(\frac{n\pi}{l}\right)^2,\qquad
 X_n(x)=\cos\frac{n\pi x}{l},\qquad n=1,2,\ldots.
\]
Нулевая мода не затухает; все остальные моды затухают экспоненциально.
Для $\phi\in L^2(0,l)$ решение имеет вид
\[
\boxed{u(x,t)=\overline\phi+
 \sum_{n=1}^{\infty}A_n
 e^{-a^2(n\pi/l)^2t}\cos\frac{n\pi x}{l},}
\]
где
\[
 A_n=\frac2l\int_0^l\phi(\xi)\cos\frac{n\pi\xi}{l}\,d\xi.
\]
Ряд задаёт решение при $t>0$, а начальное условие выполняется
в смысле сходимости в $L^2(0,l)$. При достаточной гладкости и
согласовании начальных и граничных данных получается классическое
решение вплоть до начального момента.

Для $\phi(x)=x$ вычисляем коэффициенты:
\[
\begin{aligned}
 A_n&=\frac2l\int_0^lx\cos\frac{n\pi x}{l}\,dx\\
 &=\frac{2l}{n^2\pi^2}\bigl((-1)^n-1\bigr).
\end{aligned}
\]
Следовательно, коэффициенты чётных мод равны нулю, а для нечётных
\[
 A_{2m+1}=-\frac{4l}{(2m+1)^2\pi^2}.
\]
Окончательно,
\[
\boxed{\begin{aligned}
 u(x,t)=\frac l2-\frac{4l}{\pi^2}
 \sum_{m=0}^{\infty}\frac{1}{(2m+1)^2}
 &\exp\left(-\frac{a^2(2m+1)^2\pi^2}{l^2}t\right)\\
 &\times\cos\frac{(2m+1)\pi x}{l}.
\end{aligned}}
\]
При $t\to\infty$ все слагаемые суммы затухают:
\[
\boxed{\lim_{t\to\infty}u(x,t)=\frac l2.}
\]
Более того,
\[
 \sup_{0\le x\le l}\left|u(x,t)-\frac l2\right|
 \le\frac{4l}{\pi^2}e^{-a^2\pi^2t/l^2}
 \sum_{m=0}^{\infty}\frac{1}{(2m+1)^2}
 =\frac l2e^{-a^2\pi^2t/l^2}.
\]
Значит, сходимость к ненулевому стационарному режиму равномерна
по всей длине стержня.

Следует отметить, что начальная функция $\phi(x)=x$ не удовлетворяет
условиям согласования $\phi'(0)=\phi'(l)=0$.
Поэтому не следует требовать непрерывности $u_x$ в угловых точках
$(0,0)$ и $(l,0)$. Это не препятствует существованию решения:
начальная температура восстанавливается при $t\downarrow0$, а
теплоизоляция торцов выполняется для каждого $t>0$.

\section{Необходимое и достаточное условие}

Для выбранных однородных условий Неймана и произвольной
$\phi\in L^2(0,l)$ предельный стационарный режим равен
\[
 U_\infty(x)=\overline\phi.
\]
Поэтому необходимое и достаточное условие его ненулевости имеет вид
\[
\boxed{U_\infty\not\equiv0
 \quad\Longleftrightarrow\quad
 \int_0^l\phi(x)\,dx\ne0.}
\]

Необходимость: если предел равен ненулевой постоянной $C$, то
по сохранению средней температуры
\[
 \int_0^l\phi(x)\,dx=lC\ne0.
\]
Достаточность: если интеграл начальной температуры ненулевой, то
нулевая мода равна $\overline\phi\ne0$, а все остальные моды затухают.
В частности,
\[
 \|u(\cdot,t)-\overline\phi\|_{L^2(0,l)}
 \le e^{-a^2\pi^2t/l^2}\|\phi-\overline\phi\|_{L^2(0,l)}.
\]
Сходимость также равномерна при $t\to\infty$: для любого $t_0>0$
\[
\begin{aligned}
 \sup_x|u(x,t)-\overline\phi|
 &\le e^{-a^2\pi^2(t-t_0)/l^2}
 \left(\sum_{n=1}^{\infty}A_n^2\right)^{1/2}
 \left(\sum_{n=1}^{\infty}
 e^{-2a^2(n\pi/l)^2t_0}\right)^{1/2},\quad t\ge t_0.
\end{aligned}
\]
Правая часть стремится к нулю.

Условие $\phi\not\equiv0$ само по себе недостаточно. Например,
\[
 \phi(x)=x-\frac l2
\]
является ненулевой функцией, но её среднее равно нулю; решение
стремится к нулевому стационарному режиму. Если же
$\phi\ge0$ почти всюду и $\phi$ не равна нулю почти всюду, её
интеграл положителен, и предельный режим ненулевой.

Наконец, существование ненулевых стационарных решений краевой задачи
следует отличать от достижения ненулевого предела конкретным решением.
Постоянные $C\ne0$ допустимы при любых начальных данных, но решение
с нулевым начальным средним стремится именно к $C=0$.
Режим стационарен уже при $t=0$ тогда и только тогда, когда
$\phi$ является постоянной почти всюду; для ненулевого точного
стационарного решения эта постоянная должна быть ненулевой.

\chapter{Стационарный режим при неоднородных граничных условиях}

\section{Постановка индивидуальной задачи}

Рассмотрим стержень длины $l>0$ с теплоизолированной боковой
поверхностью и без источников тепла. Начальная температура равна $x$.
Типы условий теперь задаются индивидуальным вариантом:
первый род при $x=0$ и второй род при $x=l$.

Сохраняем знаки общей постановки с изображения:
\[
 \alpha u(0,t)-\beta u_x(0,t)=\mu(t),\qquad
 \gamma u(l,t)-\delta u_x(l,t)=\nu(t).
\]
Выбираем нормировку
\[
 \alpha=1,\quad\beta=0,\qquad
 \gamma=0,\quad\delta=1,\qquad
 \mu(t)=A,\quad\nu(t)=B.
\]
Следовательно, индивидуальные неоднородные условия имеют вид
\[
\boxed{u(0,t)=A,\qquad -u_x(l,t)=B.}
\]
Знак минус во втором условии обусловлен записью исходного граничного
оператора. Поэтому здесь $B=-u_x(l,t)$, а не $B=u_x(l,t)$.
Параметры $A,B$ --- произвольные вещественные константы.

Полная задача:
\[
\boxed{\begin{cases}
 u_t=a^2u_{xx},&0<x<l,\quad t>0,\\
 u(0,t)=A,\quad -u_x(l,t)=B,&t>0,\\
 u(x,0)=x,&0\le x\le l.
\end{cases}}
\]
Теплоизоляция боковой поверхности означает $b=0$; отсутствие
распределённых источников означает $f=0$. При этом торцы могут
обмениваться теплом с внешней средой: это воздействие задают
неоднородные граничные условия.

\section{Нахождение стационарного профиля}

В стационарном режиме температура не зависит от времени:
\[
 u(x,t)=U(x),\qquad u_t=0.
\]
Следовательно,
\[
 U''(x)=0,\qquad U(0)=A,\qquad -U'(l)=B.
\]
Общее решение дифференциального уравнения имеет вид
\[
 U(x)=C_1x+C_2.
\]
Из левого граничного условия получаем $C_2=A$, а из правого
$-C_1=B$. Поэтому
\[
\boxed{U(x;A,B)=A-Bx,\qquad 0\le x\le l.}
\]
Проверка:
\[
 U''=0,\qquad U(0)=A,\qquad -U'(l)=B.
\]
Стационарный профиль существует и единственен при любых $A,B$.
Действительно, разность двух стационарных решений удовлетворяет
$W''=0$, $W(0)=0$, $W'(l)=0$, откуда $W\equiv0$.

Если условие второго рода нормируют иначе, записывая $u_x(l,t)=B$,
ответ имеет вид $U=A+Bx$. Это другой выбор знака параметра $B$.
В настоящем разделе используется именно форма $-u_x(l,t)=B$
из исходной постановки.

\section{Установление стационарного режима}

Чтобы показать, что найденный профиль является предельным режимом,
выполним замену
\[
 u(x,t)=U(x)+v(x,t)=A-Bx+v(x,t).
\]
Тогда
\[
\begin{cases}
 v_t=a^2v_{xx},&0<x<l,\ t>0,\\
 v(0,t)=0,\quad v_x(l,t)=0,&t>0,\\
 v(x,0)=(1+B)x-A.
\end{cases}
\]
Собственные значения и функции однородной смешанной задачи равны
\[
 k_n=\frac{(n+\tfrac12)\pi}{l},\qquad
 \lambda_n=k_n^2,\qquad X_n(x)=\sin(k_nx),\qquad n=0,1,2,\ldots.
\]
Поэтому решение записывается в виде
\[
 u(x,t)=A-Bx+
 \sum_{n=0}^{\infty}c_n e^{-a^2k_n^2t}\sin(k_nx),
\]
где
\[
 c_n=\frac2l\int_0^l\bigl((1+B)\xi-A\bigr)\sin(k_n\xi)\,d\xi.
\]
Используя $\cos(k_nl)=0$ и $\sin(k_nl)=(-1)^n$, получаем
\[
 \int_0^l\sin(k_n\xi)\,d\xi=\frac1{k_n},\qquad
 \int_0^l\xi\sin(k_n\xi)\,d\xi=\frac{(-1)^n}{k_n^2}.
\]
Таким образом,
\[
\boxed{c_n=\frac2l\left(
 \frac{(1+B)(-1)^n}{k_n^2}-\frac A{k_n}\right).}
\]
Все собственные значения положительны, и переходные моды затухают:
\[
\boxed{u(\cdot,t)\longrightarrow A-Bx\quad(t\to\infty).}
\]
В частности, в норме $L^2(0,l)$
\[
 \|u(\cdot,t)-U\|_{L^2}
 \le e^{-a^2\pi^2t/(4l^2)}\|x-U\|_{L^2}.
\]
Для произвольного начального распределения $\phi\in L^2(0,l)$
формула сохраняется с коэффициентами
\[
 c_n=\frac2l\int_0^l\bigl(\phi(\xi)-A+B\xi\bigr)
 \sin(k_n\xi)\,d\xi.
\]
Следовательно, предельный профиль не зависит от начальных условий.

Начальное распределение $x$ согласовано с обоими граничными условиями
в начальный момент только при $A=0$, $B=-1$.
Для остальных значений решение рассматривается при $t>0$ с
начальным условием в смысле $L^2$; гладкость вплоть до угловых точек
$(0,0)$ и $(l,0)$ не предполагается.
При $A=0$, $B=-1$ начальная функция совпадает со стационарным
профилем, и $u(x,t)=x$ для всех $t$.

\section{Зависимость от параметров}

Профиль линейно зависит от $A$ и $B$:
\[
 U(x;A,B)=A-Bx,\qquad
 \frac{\partial U}{\partial A}=1,\qquad
 \frac{\partial U}{\partial B}=-x.
\]
Изменение $A$ на $\Delta A$ сдвигает весь профиль на одну и ту же
величину. Изменение $B$ на $\Delta B$ меняет температуру на $-x\Delta B$:
влияние отсутствует на левом торце и максимально по модулю на правом.
Точное приращение равно
\[
 \Delta U(x)=\Delta A-x\Delta B,
 \qquad
 \|\Delta U\|_{\infty}\le|\Delta A|+l|\Delta B|.
\]

Температуры концов и средняя температура:
\[
 U(0)=A,\qquad U(l)=A-Bl,\qquad
 \overline U=\frac1l\int_0^lU(x)\,dx=A-\frac{Bl}{2}.
\]
Если $B>0$, температура убывает вдоль стержня;
если $B<0$, возрастает; если $B=0$, профиль постоянен и равен $A$.
Максимальное и минимальное значения равны
\[
 \max U=\max\{A,A-Bl\},\qquad
 \min U=\min\{A,A-Bl\}.
\]
Неотрицательность температуры на всём отрезке эквивалентна условиям
\[
 U(x)\ge0\ \text{для всех }x\in[0,l]
 \quad\Longleftrightarrow\quad A\ge0,\quad A-Bl\ge0.
\]
Строгая положительность на замкнутом отрезке эквивалентна двум
строгим неравенствам. Эти ограничения имеют физический смысл,
если температура понимается как абсолютная; для температуры
относительно выбранного отсчёта отрицательные значения допустимы.

Ненулевой стационарный профиль получается тогда и только тогда, когда
\[
\boxed{U\not\equiv0\quad\Longleftrightarrow\quad (A,B)\ne(0,0).}
\]
Нулевое среднее $A=Bl/2$ не означает тождественного нуля:
при $B\ne0$ профиль меняет знак и остаётся ненулевым.
При $B\ne0$ его нуль расположен в точке $x_*=A/B$,
если эта точка принадлежит отрезку $[0,l]$.

\section{Физическая интерпретация}

Согласно закону Фурье плотность потока в направлении положительной
оси $Ox$ равна
\[
 j_x=-kU'=kB.
\]
Поток через поперечное сечение площадью $S$ равен $Q=kSB$.
Таким образом, $B$ задаёт поток с точностью до множителя $kS$,
а $A$ задаёт температуру левого торца.
При $B>0$ поток направлен вправо; при $B<0$ --- влево.

Внешняя нормаль на левом торце направлена в сторону $-x$,
а на правом --- в сторону $+x$. Поэтому наружные потоки равны
\[
 Q_{\mathrm{out},0}=-kSB,\qquad
 Q_{\mathrm{out},l}=kSB,
 \qquad Q_{\mathrm{out},0}+Q_{\mathrm{out},l}=0.
\]
В стационарном режиме тепло не накапливается: поток, поступающий
через один торец, выходит через другой. Ненулевой градиент температуры
и ненулевой поток совместимы со стационарностью.
Коэффициент $a^2$ не входит в формулу $U=A-Bx$, но определяет скорость
затухания переходного процесса.

\chapter{Неоднородные условия на торце и асимптотика решения}

\section{Постановка задачи}

Сохраняем индивидуальные типы условий: первый род при $x=0$,
второй род при $x=l$. Рассмотрим первую задачу в общем виде
\[
\begin{cases}
 u_t=a^2u_{xx}-bu+f(x,t),&0<x<l,\quad t>0,\\
 u(0,t)=g(t)=A\cos(\omega_1t)+B\sin(\omega_2t),\\
 -u_x(l,t)=0,\\
 u(x,0)=\phi(x).
\end{cases}
\]
Здесь $a>0$, $b\ge0$, $A,B\in\mathbb R$, $\omega_1,\omega_2>0$.
Поскольку правое условие однородно, оно эквивалентно $u_x(l,t)=0$.

В отличие от предыдущей задачи, первая задача допускает боковой
теплообмен ($b>0$) и распределённый источник $f$.
Поэтому ниже сначала приводится решение для произвольных $f,\phi$.
Исследование воздействия только граничной неоднородности выполняется
при $f=0$; иначе асимптотику определяет также источник.
Для теплоизолированной боковой поверхности дополнительно полагаем
$b=0$. Начальную температуру $\phi(x)=x$ можно подставить во все формулы.

\section{Замена и метод Фурье}

Положим
\[
 u(x,t)=g(t)+v(x,t).
\]
Тогда
\[
\begin{cases}
 v_t=a^2v_{xx}-bv+F(x,t),\\
 v(0,t)=0,\quad v_x(l,t)=0,\\
 v(x,0)=\phi(x)-g(0),
\end{cases}
\qquad F=f-g'-bg.
\]
Собственные функции и скорости затухания:
\[
 X_n(x)=\sin(k_nx),\quad k_n=\frac{(n+\tfrac12)\pi}{l},\quad
 \sigma_n=a^2k_n^2+b>0,\qquad n=0,1,\ldots.
\]
Ищем $v=\sum_{n=0}^{\infty}v_n(t)X_n(x)$.
Обозначим
\[
 d_n=\frac2l\int_0^lX_n(x)\,dx=\frac{2}{lk_n},\qquad
 f_n(t)=\frac2l\int_0^lf(\xi,t)\sin(k_n\xi)\,d\xi,
\]
\[
 c_n=\frac2l\int_0^l\bigl(\phi(\xi)-g(0)\bigr)\sin(k_n\xi)\,d\xi.
\]
Получаем
\[
 v_n'+\sigma_nv_n=f_n-d_n(g'+bg),\qquad v_n(0)=c_n.
\]
Следовательно, решение исходной задачи равно
\[
\boxed{\begin{aligned}
 u(x,t)=g(t)+\sum_{n=0}^{\infty}\sin(k_nx)
 \bigg[&c_ne^{-\sigma_nt}\\
 &+\int_0^te^{-\sigma_n(t-\tau)}
 \bigl(f_n(\tau)-d_n(g'(\tau)+bg(\tau))\bigr)\,d\tau\bigg].
\end{aligned}}
\]
Для заданной гармонической неоднородности
\[
 g'+bg=A\bigl(b\cos\omega_1t-\omega_1\sin\omega_1t\bigr)
 +B\bigl(b\sin\omega_2t+\omega_2\cos\omega_2t\bigr).
\]
Если $\phi(x)=x$, то $g(0)=A$ и
\[
 c_n=\frac2l\left(\frac{(-1)^n}{k_n^2}-\frac A{k_n}\right).
\]
Несогласованность начальных данных с граничными значениями в углах
не препятствует решению при $t>0$; начальное условие можно понимать
в смысле $L^2(0,l)$.

\section{Гармонический установившийся режим}

Пусть $f=0$. Введём комплексную функцию
\[
 q_\omega=\frac{\sqrt{b+i\omega}}{a},\qquad
 H_\omega(x)=\frac{\cosh(q_\omega(l-x))}{\cosh(q_\omega l)}.
\]
Берётся ветвь корня с положительной действительной частью.
Функция удовлетворяет условиям
\[
 a^2H_\omega''-(b+i\omega)H_\omega=0,\qquad
 H_\omega(0)=1,\quad H_\omega'(l)=0.
\]
Поэтому точное частное решение имеет вид
\[
\boxed{P(x,t)=A\operatorname{Re}\bigl(H_{\omega_1}(x)e^{i\omega_1t}\bigr)
 +B\operatorname{Im}\bigl(H_{\omega_2}(x)e^{i\omega_2t}\bigr).}
\]
Оно удовлетворяет уравнению, условию $P(0,t)=g(t)$ и $P_x(l,t)=0$.
Полное решение:
\[
\boxed{u(x,t)=P(x,t)+\sum_{n=0}^{\infty}r_n
 e^{-\sigma_nt}\sin(k_nx),}
\]
\[
 r_n=\frac2l\int_0^l\bigl(\phi(\xi)-P(\xi,0)\bigr)
 \sin(k_n\xi)\,d\xi.
\]
Таким образом,
\[
 \|u(\cdot,t)-P(\cdot,t)\|_{L^2}
 \le e^{-\sigma_0t}\|\phi-P(\cdot,0)\|_{L^2},\qquad
 \sigma_0=b+\frac{a^2\pi^2}{4l^2}.
\]
После любого положительного времени сглаживание также даёт
равномерное затухание переходной части.

При ненулевых амплитудах и положительных частотах решение в общем
случае не имеет предела: оно приближается к ограниченному колебательному
режиму. Если активна одна частота или отношение двух активных частот
рационально, режим периодический. При двух ненулевых амплитудах и
иррациональном отношении частот он квазипериодический.
Начальные условия влияют только на затухающую часть.

Следовательно, при фиксированной форме $A\cos\omega_1t+B\sin\omega_2t$
и $f=0$ неограниченный режим невозможен. При положительных частотах
ненулевой постоянный предел также невозможен. Выбор $A=B=0$
даёт затухание к нулю. Если допускается $\omega_1=0$, косинусный член
становится постоянным и возникает стационарный профиль, рассмотренный ниже.
Чтобы получить все требуемые типы поведения, далее подбираем другие
временные зависимости $g(t)$, сохраняя типы границ и правое условие нулевым.

\section{Аналитический подбор воздействий}

Во всех трёх примерах этой части считаем $f=0$, а начальную функцию
$\phi$ произвольной, в частности $\phi(x)=x$.

\subsection{Решение стремится к нулю}

Выберем ненулевое затухающее воздействие
\[
 g(t)=C e^{-rt},\qquad C\ne0,\qquad 0<r<\sigma_0.
\]
Определим $H_r$ из
\[
 a^2H_r''-(b-r)H_r=0,\qquad H_r(0)=1,\quad H_r'(l)=0.
\]
Явно,
\[
 H_r(x)=\begin{cases}
 \displaystyle\frac{\cosh(\kappa(l-x))}{\cosh(\kappa l)},
 &b>r,\quad\kappa=\sqrt{b-r}/a,\\[0.7em]
 1,&b=r,\\[0.4em]
 \displaystyle\frac{\cos(\eta(l-x))}{\cos(\eta l)},
 &b<r,\quad\eta=\sqrt{r-b}/a.
 \end{cases}
\]
Условие $r<\sigma_0$ исключает нуль знаменателя.
Точное частное решение $P_r=C e^{-rt}H_r$ даёт
\[
 u(x,t)=C e^{-rt}H_r(x)+\sum_{n=0}^{\infty}z_n
 e^{-\sigma_nt}\sin(k_nx),
\]
где $z_n$ --- коэффициенты Фурье функции $\phi-CH_r$.
Обе части затухают, поэтому
\[
\boxed{u(x,t)\longrightarrow0.}
\]

\subsection{Решение неограниченно возрастает}

Выберем
\[
 g(t)=C e^{rt},\qquad C>0,\quad r>0.
\]
Пусть
\[
 \kappa_r=\frac{\sqrt{b+r}}{a},\qquad
 K_r(x)=\frac{\cosh(\kappa_r(l-x))}{\cosh(\kappa_rl)}>0.
\]
Тогда $P_r=C e^{rt}K_r$ --- точное частное решение, и
\[
 u(x,t)=C e^{rt}K_r(x)+\sum_{n=0}^{\infty}z_n
 e^{-\sigma_nt}\sin(k_nx),
\]
где $z_n$ --- коэффициенты функции $\phi-CK_r$.
Поскольку $K_r$ имеет положительный минимум на $[0,l]$,
\[
\boxed{u(x,t)\longrightarrow+\infty}
\]
в каждой точке и равномерно по $x$. Неограниченность следует не только
из роста температуры на границе, но и из явного решения внутри стержня.

\subsection{Постоянный предельный режим}

Выберем $g(t)=C$, $C\ne0$. При $b=0$ точное частное решение равно $C$:
\[
 u(x,t)=C+\sum_{n=0}^{\infty}z_n
 e^{-a^2k_n^2t}\sin(k_nx),\qquad
 z_n=\frac2l\int_0^l(\phi(\xi)-C)\sin(k_n\xi)\,d\xi.
\]
Таким образом, для теплоизолированной боковой поверхности
\[
\boxed{u(x,t)\longrightarrow C.}
\]

При $b>0$ и $f=0$ то же граничное воздействие приводит не к
пространственной константе, а к стационарному профилю
\[
\boxed{U_C(x)=C\frac{\cosh(\sqrt b(l-x)/a)}{\cosh(\sqrt bl/a)}.}
\]
Переходная часть вновь является суммой затухающих мод.
Важно различать постоянство во времени и постоянство по координате.
Если требуется именно $U\equiv C\ne0$, то при $b>0$, $f=0$
такой профиль невозможен при любом выборе граничного воздействия:
подстановка в стационарное уравнение даёт $0=-bC$.
Для его поддержания нужен источник $f=bC$ и граничные данные
$u(0,t)=C$, $u_x(l,t)=0$. При этих изменённых данных $u-C$
решает однородную затухающую задачу.

\section{Роль распределённого источника}

При произвольном $f(x,t)$ поведение нельзя определить только по $g$.
К решению с $f=0$ добавляется
\[
 u_f(x,t)=\sum_{n=0}^{\infty}\sin(k_nx)
 \int_0^te^{-\sigma_n(t-\tau)}f_n(\tau)\,d\tau.
\]
В частности, постоянный источник обычно создаёт ненулевой стационарный
вклад, даже если граничное воздействие затухает; растущий источник
может создавать неограниченный режим при нулевых границах.
Поэтому приведённый аналитический подбор изолирует действие граничных
неоднородностей условием $f=0$ и явно различает случаи $b=0$ и $b>0$.

\chapter{Боковой теплообмен по закону Ньютона}

\section{Постановка практического задания 4}

В пункте 4 практической части требуется решить вторую практическую
задачу с учётом теплообмена на боковой поверхности. Поэтому сохраняются
начальная температура $u(x,0)=x$ и постоянные индивидуальные
неоднородности $A,B$ на торцах. Гармонические данные из практического
задания 3 здесь не используются.

Для варианта 11 типы условий: первый род при $x=0$ и второй род
при $x=l$. Сохраняя знак общей записи лабораторной работы,
получаем
\[
 u(0,t)=A,\qquad -u_x(l,t)=B.
\]
Здесь $B$ обозначает отрицательную координатную производную
температуры. Если вместо этого определить $B=u_x(l,t)$, во всех
последующих формулах следует заменить $B$ на $-B$.

Обозначим постоянную температуру среды через $T_e$.
Для однородного стержня с площадью сечения $S$ и периметром сечения $P$
закон Ньютона на боковой поверхности даёт уравнение
\[
 cS u_t=kS u_{xx}-HP(u-T_e).
\]
Здесь $c$ --- объёмная теплоёмкость, $k$ --- теплопроводность,
$H>0$ --- коэффициент теплоотдачи. Внутренние источники отсутствуют.
Введём
\[
 a^2=\frac{k}{c},\qquad b=\frac{HP}{cS}>0.
\]
Для круглого сечения радиуса $R$ имеем $b=2H/(cR)$.
Таким образом,
\[
\boxed{\begin{cases}
 u_t=a^2u_{xx}-b(u-T_e),&0<x<l,\quad t>0,\\
 u(0,t)=A,\quad -u_x(l,t)=B,&t>0,\\
 u(x,0)=x,&0\le x\le l.
\end{cases}}
\]
В обозначениях $u_t=a^2u_{xx}-bu+f$ функция $f=bT_e$ возникает
из теплообмена, а не из дополнительного внутреннего источника.

\section{Стационарный профиль при постоянной среде}

Пусть $u(x,t)=U_{T_e}(x)$. Тогда
\[
 a^2U_{T_e}''-b(U_{T_e}-T_e)=0,\qquad
 U_{T_e}(0)=A,\quad U_{T_e}'(l)=-B.
\]
Обозначим
\[
 \kappa=\frac{\sqrt b}{a}>0.
\]
Положим $Y=U_{T_e}-T_e$. Тогда
\[
 Y''-\kappa^2Y=0,\qquad Y(0)=A-T_e,\quad Y'(l)=-B.
\]
Удобно представить общее решение относительно правого торца:
\[
 Y(x)=C\cosh(\kappa(l-x))+D\sinh(\kappa(l-x)).
\]
Из $Y'(l)=-B$ получаем $D=B/\kappa$, а из левого условия
\[
 C=\frac{A-T_e-(B/\kappa)\sinh(\kappa l)}{\cosh(\kappa l)}.
\]
Применяя формулу разности гиперболических синусов, получаем
\[
\boxed{U_{T_e}(x)=T_e+
 (A-T_e)\frac{\cosh(\kappa(l-x))}{\cosh(\kappa l)}
 -\frac{B}{\kappa}\frac{\sinh(\kappa x)}{\cosh(\kappa l)}.}
\]
Непосредственная проверка даёт
\[
 U_{T_e}(0)=A,\qquad U_{T_e}'(l)=-B,\qquad
 U_{T_e}''=\kappa^2(U_{T_e}-T_e).
\]
Стационарная задача имеет единственное решение. Действительно,
для разности $W$ двух решений имеем $a^2W''-bW=0$,
$W(0)=0$, $W'(l)=0$. Умножая уравнение на $W$ и интегрируя,
получаем
\[
 -a^2\int_0^l(W')^2\,dx-b\int_0^lW^2\,dx=0,
\]
откуда $W\equiv0$.

\section{Полное решение методом Фурье}

Представим температуру как сумму стационарной и переходной частей:
\[
 u(x,t)=U_{T_e}(x)+v(x,t).
\]
Тогда
\[
\begin{cases}
 v_t=a^2v_{xx}-bv,\\
 v(0,t)=0,\quad v_x(l,t)=0,\\
 v(x,0)=x-U_{T_e}(x).
\end{cases}
\]
Собственные функции и скорости затухания:
\[
 X_n(x)=\sin(k_nx),\quad
 k_n=\frac{(n+\tfrac12)\pi}{l},\quad
 \sigma_n=a^2k_n^2+b,\qquad n=0,1,2,\ldots.
\]
Следовательно,
\[
\boxed{u(x,t)=U_{T_e}(x)+\sum_{n=0}^{\infty}
 C_n(T_e)e^{-\sigma_nt}\sin(k_nx),}
\]
где
\[
\boxed{C_n(T_e)=\frac2l\int_0^l
 \bigl(\xi-U_{T_e}(\xi)\bigr)\sin(k_n\xi)\,d\xi.}
\]
Это уже явная формула решения. Можно вычислить коэффициенты
без интегралов. Поскольку $\cos(k_nl)=0$, $\sin(k_nl)=(-1)^n$,
\[
 \int_0^l\xi\sin(k_n\xi)\,d\xi=\frac{(-1)^n}{k_n^2}.
\]
Умножим равенство $U_{T_e}''=\kappa^2(U_{T_e}-T_e)$
на $\sin(k_nx)$ и дважды проинтегрируем по частям.
Граничные члены дают
\[
 (k_n^2+\kappa^2)\int_0^lU_{T_e}(x)\sin(k_nx)\,dx
 =Ak_n-B(-1)^n+\frac{\kappa^2T_e}{k_n}.
\]
Таким образом,
\[
\boxed{C_n(T_e)=\frac2l\left[
 \frac{(-1)^n}{k_n^2}
 -\frac{Ak_n-B(-1)^n+\kappa^2T_e/k_n}{k_n^2+\kappa^2}
 \right].}
\]
Для произвольной начальной функции $\phi$ вместо $x$
в интегральной формуле используется $\phi(\xi)$.
При $\phi=x$ согласование с границами в момент $t=0$
возможно лишь при $A=0$, $B=-1$.
Для остальных значений начальное условие понимается в смысле
$L^2(0,l)$, а решение гладко при $t>0$.
Даже при согласовании $\phi=x$ не является стационарной
температурой при $b>0$ и постоянном $T_e$.

\section{Два случая температуры среды}

\subsection{Температура среды равна нулю}

При $T_e=0$ уравнение имеет вид
\[
 u_t=a^2u_{xx}-bu.
\]
Стационарная температура равна
\[
\boxed{U_0(x)=A\frac{\cosh(\kappa(l-x))}{\cosh(\kappa l)}
 -\frac B\kappa\frac{\sinh(\kappa x)}{\cosh(\kappa l)}.}
\]
Полное решение:
\[
 u(x,t)=U_0(x)+\sum_{n=0}^{\infty}
 C_n(0)e^{-(a^2k_n^2+b)t}\sin(k_nx),
\]
\[
 C_n(0)=\frac2l\left[
 \frac{(-1)^n}{k_n^2}
 -\frac{Ak_n-B(-1)^n}{k_n^2+\kappa^2}\right].
\]
Если $A=B=0$, то $U_0\equiv0$ и
\[
 u(x,t)=\frac2l\sum_{n=0}^{\infty}\frac{(-1)^n}{k_n^2}
 e^{-(a^2k_n^2+b)t}\sin(k_nx)\longrightarrow0.
\]
Если хотя бы одна граничная неоднородность ненулевая,
стационарный профиль ненулевой: торцы поддерживают температуру
или поток, компенсирующие боковые потери.

\subsection{Температура среды постоянна}

При $T_e=T_0=\mathrm{const}$ получаем
\[
 u_t=a^2u_{xx}-bu+bT_0.
\]
Стационарная температура и полное решение:
\[
\boxed{U_{T_0}(x)=T_0+
 (A-T_0)\frac{\cosh(\kappa(l-x))}{\cosh(\kappa l)}
 -\frac B\kappa\frac{\sinh(\kappa x)}{\cosh(\kappa l)},}
\]
\[
 u(x,t)=U_{T_0}(x)+\sum_{n=0}^{\infty}
 C_n(T_0)e^{-(a^2k_n^2+b)t}\sin(k_nx).
\]
При одинаковых $A,B$ отличие от случая нулевой среды равно
\[
\boxed{U_{T_0}(x)-U_0(x)=T_0\left[
 1-\frac{\cosh(\kappa(l-x))}{\cosh(\kappa l)}\right].}
\]
На левом торце разность нулевая: его температура фиксирована
и равна $A$. Поэтому нельзя просто прибавить $T_0$ к $U_0$,
не изменив одновременно левое граничное значение.

Если $A=T_0$ и $B=0$, то
\[
 U_{T_0}(x)\equiv T_0.
\]
Это необходимое и достаточное условие пространственно постоянного
стационарного профиля: постоянство требует $U=T_0$
из уравнения и $A=T_0$, $B=0$ из граничных условий.
Если $A=B=0$, но $T_0\ne0$, профиль ненулевой:
\[
 U_{T_0}(x)=T_0\left[
 1-\frac{\cosh(\kappa(l-x))}{\cosh(\kappa l)}\right].
\]
При $b>0$ тождественно нулевой стационарный режим возможен
тогда и только тогда, когда $A=B=T_0=0$.

\section{Асимптотика и проверка предела}

Все скорости затухания положительны:
\[
 \sigma_n=a^2k_n^2+b>0.
\]
Поэтому в обоих случаях
\[
\boxed{u(\cdot,t)\longrightarrow U_{T_e}\quad(t\to\infty).}
\]
Начальное распределение влияет только на переходный процесс.
В норме $L^2(0,l)$ справедлива оценка
\[
 \|u(\cdot,t)-U_{T_e}\|_{L^2}
 \le e^{-(b+a^2\pi^2/(4l^2))t}
 \|x-U_{T_e}\|_{L^2}.
\]
Сходимость равномерна по $x$ при $t\to\infty$ вследствие
сглаживания решения и экспоненциального затухания ряда.
Боковой теплообмен увеличивает скорости затухания на $b$.
При этом стационарный профиль зависит от $b$ через $\kappa$.

Полезная проверка: при исчезновении теплообмена, $b\downarrow0$,
\[
 \frac{\cosh(\kappa(l-x))}{\cosh(\kappa l)}\longrightarrow1,
 \qquad
 \frac{\sinh(\kappa x)}{\kappa\cosh(\kappa l)}\longrightarrow x.
\]
Поэтому
\[
\boxed{U_{T_e}(x)\longrightarrow A-Bx.}
\]
Температура среды исчезает из предельной формулы, как и должно быть
при отсутствии бокового теплообмена. Это совпадает со стационарным
режимом второй практической задачи.

\chapter{Практическое задание 5. Нагрев шара постоянным потоком}

\section{Постановка и знак теплового потока}

Однородный шар радиуса $R$ имеет нулевую начальную температуру.
На его поверхность с момента $t=0$ равномерно подаётся постоянный
поток тепла. Обозначим через $q>0$ плотность входящего теплового потока
(теплота на единицу площади за единицу времени), через $k>0$ ---
теплопроводность, через $c>0$ --- объёмную теплоёмкость.
Положим $a^2=k/c$.
Если в условии $q$ означает полный поток через всю поверхность,
в формулах вместо $q$ следует использовать $q/(4\pi R^2)$.

Из сферической симметрии температура зависит только от радиуса $r$ и
времени. Уравнение и условия имеют вид
\[
\boxed{\begin{cases}
 u_t=a^2\left(u_{rr}+\dfrac2r u_r\right),&0<r<R,\ t>0,\\
 u_r(0,t)=0,\quad u(0,t)\text{ конечна},\\
 k u_r(R,t)=q,\\
 u(r,0)=0.
\end{cases}}
\]
Наружная плотность потока по закону Фурье равна $-ku_r(R,t)$.
Поскольку шар нагревается, наружный поток равен $-q$, откуда
$ku_r(R,t)=q$. Эта проверка определяет знак граничного условия.

\section{Баланс тепла и выделение растущей части}

Средняя по объёму температура равна
\[
 \overline u(t)=\frac3{R^3}\int_0^R u(r,t)r^2\,dr.
\]
Интегрирование уравнения даёт
\[
 \overline u'(t)=\frac{3a^2}{R^3}[r^2u_r]_0^R
 =\frac{3q}{cR},\qquad
 \boxed{\overline u(t)=\frac{3q}{cR}t.}
\]
Следовательно, при $q>0$ стационарная температура невозможна:
тепло непрерывно поступает в шар и не отводится.

Выделим точное частное решение
\[
 p(r,t)=\frac{3q}{cR}t+
 \frac{qR}{k}\left(\frac{r^2}{2R^2}-\frac3{10}\right).
\]
Оно удовлетворяет уравнению и условию $kp_r(R,t)=q$.
Постоянная $-3/10$ выбрана так, чтобы среднее пространственной
части было нулевым. Положим $u=p+v$. Тогда
\[
 v_t=a^2\left(v_{rr}+\frac2r v_r\right),\qquad
 v_r(0,t)=v_r(R,t)=0,
\]
\[
 v(r,0)=-\frac{qR}{k}\left(\frac{r^2}{2R^2}-\frac3{10}\right).
\]
Начальное среднее $v$ равно нулю; постоянная собственная мода отсутствует.

\section{Собственные функции и коэффициенты}

Разделение переменных приводит к задаче
\[
 (r^2X')'+\lambda r^2X=0,\qquad X\text{ регулярна при }r=0,
 \qquad X'(R)=0.
\]
Для положительных собственных значений удобно взять
\[
 X_n(r)=\frac{\sin(z_nr/R)}{z_nr/R},\qquad X_n(0)=1,
 \qquad\lambda_n=\frac{z_n^2}{R^2}.
\]
Условие при $r=R$ даёт
\[
 z_n\cos z_n-\sin z_n=0,
 \qquad \boxed{\tan z_n=z_n.}
\]
Здесь $z_n$ --- строго положительные ненулевые корни:
\[
 z_1\approx4.493409,\qquad z_2\approx7.725252,\qquad
 z_3\approx10.904122.
\]
Корень $z=0$ соответствует постоянной моде и рассматривается отдельно.
Ортогональность осуществляется с весом $r^2$.
Поэтому
\[
 v(r,t)=\sum_{n=1}^{\infty}D_nX_n(r)e^{-a^2z_n^2t/R^2},
 \qquad
 D_n=\frac{\int_0^R v(r,0)X_n(r)r^2\,dr}
 {\int_0^RX_n^2(r)r^2\,dr}.
\]
Положив $s=r/R$ и используя $\sin z_n=z_n\cos z_n$, получаем
\[
 \int_0^1s^2\frac{\sin(z_ns)}{z_ns}\,ds=0,
\]
\[
 \int_0^1s^2\left(\frac{s^2}2-\frac3{10}\right)
 \frac{\sin(z_ns)}{z_ns}\,ds=\frac{\cos z_n}{z_n^2},
\]
\[
 \int_0^1s^2\left(\frac{\sin(z_ns)}{z_ns}\right)^2ds
 =\frac{\cos^2z_n}{2}.
\]
Следовательно,
\[
 \boxed{D_n=-\frac{2qR}{kz_n^2\cos z_n}.}
\]

\section{Распределение температуры и асимптотика}

Для $0\le r\le R$, $t>0$ окончательно получаем
\[
\boxed{\begin{aligned}
 u(r,t)={}&\frac{3q}{cR}t+
 \frac{qR}{k}\left(\frac{r^2}{2R^2}-\frac3{10}\right)\\
 &-\frac{2qR}{k}\sum_{n=1}^{\infty}
 \frac{e^{-a^2z_n^2t/R^2}}{z_n^2\cos z_n}
 \frac{\sin(z_nr/R)}{z_nr/R}.
\end{aligned}}
\]
При $r=0$ отношение синуса к аргументу заменяется его пределом $1$:
\[
 u(0,t)=\frac{3q}{cR}t-\frac{3qR}{10k}
 -\frac{2qR}{k}\sum_{n=1}^{\infty}
 \frac{e^{-a^2z_n^2t/R^2}}{z_n^2\cos z_n}.
\]
На поверхности шара
\[
 u(R,t)=\frac{3q}{cR}t+\frac{qR}{5k}
 -\frac{2qR}{k}\sum_{n=1}^{\infty}\frac{e^{-a^2z_n^2t/R^2}}{z_n^2}.
\]
Ряд восстанавливает нулевую начальную температуру.
Включение ненулевого потока при $t=0$ нарушает согласование
начальных и граничных производных; граничное условие выполняется
для $t>0$.

При больших временах
\[
 u(r,t)=\frac{3q}{cR}t+
 \frac{qR}{k}\left(\frac{r^2}{2R^2}-\frac3{10}\right)
 +O\left(e^{-a^2z_1^2t/R^2}\right).
\]
Температура неограниченно растёт, но её отличие от средней стремится
к постоянному пространственному профилю. В частности,
\[
 u(R,t)-u(0,t)\longrightarrow\frac{qR}{2k}.
\]
Энергетическая проверка решения:
\[
 c\frac{4\pi R^3}{3}\overline u(t)=4\pi R^2qt.
\]
Левая часть равна накопленной теплоте, правая --- теплоте,
поступившей через поверхность шара.

\chapter{Практическое задание 6. Нагрев пластины внутренним источником}

\section{Постановка задачи}

Рассмотрим однородную неограниченную пластину толщины $2R$.
Выберем координату $x$ поперёк пластины, с началом в её средней плоскости:
$-R<x<R$. Из однородности условий вдоль граней температура зависит
только от $x$ и времени $t$.

Начальная температура нулевая. С момента $t=0$ действует постоянный
равномерный объёмный источник тепла плотности $Q$ (энергия на единицу
объёма за единицу времени). Обе грани поддерживаются при нулевой
температуре. Обозначим теплопроводность через $k>0$, объёмную
теплоёмкость через $c>0$ и положим
\[
 a^2=\frac{k}{c},\qquad F_0=\frac Qc.
\]
Тогда
\[
\boxed{\begin{cases}
 u_t=a^2u_{xx}+Q/c,&-R<x<R,\quad t>0,\\
 u(-R,t)=u(R,t)=0,\\
 u(x,0)=0.
\end{cases}}
\]
Если $Q$ уже обозначает нормированную правую часть уравнения,
вместо $Q/c$ следует писать $Q$; здесь используется физическая
объёмная плотность тепловыделения.
Решение чётно по $x$. На половине пластины можно эквивалентно
задать $u_x(0,t)=0$, $u(R,t)=0$.

\section{Стационарное распределение}

Стационарная температура удовлетворяет
\[
 kU''+Q=0,\qquad U(-R)=U(R)=0.
\]
Интегрируя и используя симметрию, получаем
\[
\boxed{U(x)=\frac Q{2k}(R^2-x^2).}
\]
При $Q>0$ максимум достигается в средней плоскости:
\[
 U(0)=\frac{QR^2}{2k}.
\]
Температура конечна: выделяющееся тепло отводится через обе грани.

\section{Решение методом Фурье}

Положим $u=U+v$. Тогда
\[
 v_t=a^2v_{xx},\qquad v(-R,t)=v(R,t)=0,
 \qquad v(x,0)=-U(x).
\]
Для чётного решения собственные функции имеют вид
\[
 X_n(x)=\cos(\eta_nx),\qquad
 \eta_n=\frac{(n+\tfrac12)\pi}{R},\qquad n=0,1,\ldots.
\]
Они удовлетворяют $X_n(\pm R)=0$ и нормировке
\[
 \int_{-R}^RX_n^2(x)\,dx=R.
\]
Нечётные собственные функции не участвуют, поскольку начальные данные
и источник чётны.
Разложим стационарный профиль:
\[
 U(x)=\sum_{n=0}^{\infty}U_n\cos(\eta_nx),\qquad
 U_n=\frac1R\int_{-R}^RU(x)\cos(\eta_nx)\,dx.
\]
Поскольку $\cos(\eta_nR)=0$, $\sin(\eta_nR)=(-1)^n$,
\[
 \int_0^R(R^2-x^2)\cos(\eta_nx)\,dx
 =\frac{2(-1)^n}{\eta_n^3}.
\]
Поэтому
\[
 \boxed{U_n=\frac{2Q(-1)^n}{kR\eta_n^3}
 =\frac{16QR^2}{k\pi^3}\frac{(-1)^n}{(2n+1)^3}.}
\]
Начальная функция $v=-U$ имеет противоположные коэффициенты;
следовательно,
\[
 v(x,t)=-\sum_{n=0}^{\infty}U_n e^{-a^2\eta_n^2t}\cos(\eta_nx).
\]
Получаем две эквивалентные формы полного решения:
\[
\boxed{\begin{aligned}
 u(x,t)={}&\frac Q{2k}(R^2-x^2)\\
 &-\frac{16QR^2}{k\pi^3}\sum_{n=0}^{\infty}
 \frac{(-1)^n}{(2n+1)^3}
 e^{-a^2(2n+1)^2\pi^2t/(4R^2)}
 \cos\frac{(2n+1)\pi x}{2R},
\end{aligned}}
\]
или
\[
\boxed{\begin{aligned}
 u(x,t)=\frac{16QR^2}{k\pi^3}\sum_{n=0}^{\infty}
 \frac{(-1)^n}{(2n+1)^3}
 &\left(1-e^{-a^2(2n+1)^2\pi^2t/(4R^2)}\right)\\
 &\times\cos\frac{(2n+1)\pi x}{2R}.
\end{aligned}}
\]
Если координата отсчитывается от левой грани, $0\le y\le2R$,
нужно заменить $x$ на $y-R$.

\section{Проверка и предельный режим}

На гранях $x=\pm R$ каждое собственное слагаемое равно нулю;
при $t=0$ во второй форме решения множители $1-e^0$ равны нулю.
Уравнение проверяется либо подстановкой первой формы, либо
решением модальных уравнений для постоянного источника.
При $t>0$ ряд допускает необходимое дифференцирование.

При $t\to\infty$ переходная часть затухает и
\[
\boxed{u(x,t)\longrightarrow\frac Q{2k}(R^2-x^2).}
\]
Сходимость равномерная. Её ведущая скорость определяется первой модой:
\[
 u(x,t)-U(x)=-\frac{16QR^2}{k\pi^3}
 e^{-a^2\pi^2t/(4R^2)}\cos\frac{\pi x}{2R}
 +O\left(e^{-9a^2\pi^2t/(4R^2)}\right).
\]
Для $Q>0$ температура неотрицательна и не убывает по времени;
её максимум по толщине достигается в средней плоскости.
Средняя стационарная температура равна
\[
 \overline U=\frac1{2R}\int_{-R}^RU(x)\,dx=\frac{QR^2}{3k}.
\]
Наружные плотности потоков на обеих гранях одинаковы:
\[
 j_{\mathrm{out},+}=-kU'(R)=QR,\qquad
 j_{\mathrm{out},-}=kU'(-R)=QR.
\]
Их сумма равна $2QR$, то есть тепловыделению в пластине на единицу
площади её грани. Это проверяет стационарный энергетический баланс.
При $Q=0$ решение тождественно нулевое.


% ============================================================
% Основной текст размещайте после этой строки.
%
% \chapter{Название главы}
% \section{Название раздела}
% Текст работы.
%
% \bibliography{references}
% ============================================================

\end{document}